\begin{pastexam}[label={px:ntu-109-ee-ds-02}]{臺大電機所 109 學年度（B）第 2 題}
（單選；每題 4 分，答錯不倒扣）

Please choose the best possible answer in each of the following questions. In questions that ask
for time complexities, please choose the tightest bound.

The worst case time complexity of constructing a binary heap from an unordered array of n items is

(A)~O(n) \qquad (B)~O(n log*n) \qquad (C)~O(n log n) \qquad (D)~O(n$^2$) \qquad (E)~O(n$^2$ log n)
\end{pastexam}

\begin{pastsoln}
答案是 (A) O(n)。

用由下而上的建堆積（資結書虛擬碼~8.3）：從最後一個內部節點 $\lfloor n/2 \rfloor$ 往前，每個節點
做一次下沉。 高度為 $h$ 的節點下沉最多走 $h$ 層，而高度為 $h$ 的節點最多 $\lceil n/2^{h+1}
\rceil$ 個（高度從葉節點的 $0$ 算起），所以總工作量
\[
\sum_{h=0}^{\lfloor \lg n \rfloor} \Bigl\lceil \frac{n}{2^{h+1}} \Bigr\rceil \cdot O(h)
= O\Bigl(n \sum_{h=0}^{\infty} \frac{h}{2^{h}}\Bigr) = O(2n) = O(n),
\]
其中 $\sum_{h \ge 0} h/2^h = 2$。 資結書定理~8.1 把這件事寫成：交換次數小於 $2n$，
建堆積是 $\Theta(n)$。 下界也是 $\Omega(n)$：每個元素至少要看一次。
所以最緊的界是 $O(n)$。

(B)、(C)、(D)、(E) 都是正確的上界，但不是最緊的。 $O(n \log n)$ 是「一個一個插入」的建法，
每次插入最多上浮 $\lg n$ 層；題目要最緊的界，而由下而上的建法做得到 $O(n)$。
\end{pastsoln}
